% ============================================================
\chapter{Introduction}
\label{chap:intro}
% ============================================================

\section{Context and motivation}
\label{sec:motivation}

Portfolio optimisation allocates capital across assets to maximise return
at a given level of risk. Since Markowitz's seminal mean--variance model
\cite{markowitz1952portfolio}, the covariance (or correlation) 
matrix of returns has been the main statistical object used for 
portfolio optimisation. 
Hierarchical Risk Parity (HRP) is a recent optimisation method
\cite{lopezdeprado2016hrp} which uses the correlation matrix
to cluster assets based on their correlational distance,
producing a dendrogram (binary tree). 
Allocation then proceeds by recursive bisection.

Correlation brings two distinct deficits to that step. The first is a deficit
of \emph{information}: correlation is undirected (it records that two
assets co-move, but not which moves the other) and it can be spurious,
because two assets responding to a common cause are statistically
indistinguishable from two assets that are causally linked
\cite{simon1954spurious,
lopezdeprado2023causal}. The second is a deficit of \emph{interface}:
hierarchical clustering accepts only symmetric matrices, so even if a
directed description of the market were available, HRP and its descendants
could not represent it. If asset $A$ drives asset $B$ but not conversely,
neither the correlation matrix nor the dendrogram built from it can say so.
These deficits are separable, and this thesis treats them separately: the
first asks whether a better \emph{input} to hierarchical allocation exists;
the second, whether an allocator can be built that \emph{consumes} what the
better input contains. They become Requirements 1 and 2 of
Section~\ref{sec:question}.

Causal discovery supplies the candidate input. Score-based structure
learning (NOTEARS \cite{zheng2018notears} and its temporal extension
DYNOTEARS \cite{pamfil2020dynotears}) makes it tractable to fit, at every
monthly rebalance, a directed acyclic graph (DAG) over a hundred assets: an
estimate not only of which assets are linked but of the direction of
influence between them. A small literature now applies such graphs to equity markets
\cite{howard2025causal,tang2024trading}, but it stops short of the allocation
step: no published work sets hierarchical portfolio weights from a discovered
graph, and none asks what the graph's directionality is worth once it gets
there.

The question separates naturally, because a DAG carries two separable kinds
of information. Its \emph{skeleton}, the undirected pattern of which asset
pairs are causally linked, can be handed to an unmodified hierarchical
allocator by symmetrising the adjacency matrix. Its \emph{edge
orientations}, the direction of each link, cannot: a symmetric clustering
step discards them entirely, so exploiting them requires new allocation
machinery.
The contribution of this thesis is that machinery: the orientation-aware
allocator family, and the fixed-graph ablation design that hands every
allocator the identical graph so that any difference is attributable to
what each one reads from it. The decomposition measured in
Chapter~\ref{chap:evaluation} is what the instrument found. Either outcome
is informative. If orientation adds nothing, the skeleton was all that
mattered and cheap symmetrisation suffices; if it adds value, every
correlation-fed hierarchical allocator has been discarding available
information.

The answer matters to three audiences. For the practitioner, the research
question is one about costs: symmetrising a graph is free, exploiting orientation
requires the new allocation machinery of this thesis, and discovering the
graph at all costs hours of compute per backtest, so the decomposition
indicates which layers of causal structure justify their expense. For
the causal-finance literature, which now builds asset graphs at scale but
has not yet set portfolio weights from them, it says where in the graph
allocative value resides. And for research practice in a field
with a documented backtest-overfitting problem
\cite{bailey2014deflated,lopezdeprado2023causal}, the project demonstrates
that an effect of this size can be measured with a controlled ablation
design, a deterministic pipeline and multiplicity-corrected inference, and
reported at its estimated size whether or not it clears conventional
significance.

\section{The scientific question}
\label{sec:question}

The thesis addresses a single question, referred to throughout the report as
the research question:

\begin{quote}
\noindent\fbox{\parbox{\dimexpr\linewidth-2\fboxsep-2\fboxrule\relax}{%
\emph{When the correlation matrix in hierarchical portfolio
allocation is replaced by a causal graph, is there a performance
improvement, and how much of any improvement is attributable to the
graph's \textbf{skeleton} versus its \textbf{edge orientations}?}}}
\end{quote}

\noindent Every term is operationalised once, here, and held fixed for the
whole study:
\begin{itemize}[itemsep=1pt, topsep=3pt]
  \item \emph{Hierarchical portfolio allocation}: HRP's two-stage scheme
    (a symmetric distance matrix determines an asset ordering, and recursive
    bisection with inverse-cluster-variance splits determines the weights),
    long-only, fully invested, on the ${\sim}100$ largest S\&P~500
    constituents (a fixed 99-name universe), rebalanced monthly over
    2007--2024 (215 rebalances at windows up to one year; the 18-month and
    two-year windows' longer burn-ins leave 209 and 203).
  \item \emph{The correlation matrix}: the graph-blind control \HRP{} is
    López de Prado's original construction, with the distance
    $\sqrt{\tfrac{1}{2}(1-\rho_{ij})}$ on the lookback window's sample
    correlation \cite{lopezdeprado2016hrp}.
  \item \emph{A discovered causal graph}: the contemporaneous asset--asset
    block of a DYNOTEARS fit on the joint driver--asset panel at each
    rebalance (primary), with two controls on its orientation content:
    VARLiNGAM \cite{hyvarinen2010varlingam}, the same construction with
    differently-identified orientations, and a ridge-regularised Granger
    graph \cite{granger1969investigating}, which preserves a directional
    ordering but degrades the edge magnitudes. Four estimation windows are
    used (189, 252, 378 and 504 trading days; nine months to two years), so
    the decomposition is measured as a function of the estimation window
    rather than at a single conventional one.
  \item \emph{Performance improvement}: out-of-sample (walk-forward), net-of-cost
    (5\,bps per unit one-way turnover, swept over $\{0,5,10,20\}$\,bps),
    annualised Sharpe ratio of daily portfolio returns, supplemented by the
    Probabilistic and Deflated Sharpe ratios
    \cite{bailey2012psr,bailey2014deflated} because daily returns here have
    excess kurtosis ${\approx}\,\rsKurtosis$, and adjudicated family-wise by
    White's Reality Check, Hansen's SPA and the Model Confidence Set
    \cite{white2000reality,hansen2005spa,hansen2011mcs}; the Deflated
    Sharpe ratio deflates against all \rsNtrials\ configurations evaluated
    in the project, while the Reality Check, SPA and MCS each run over the
    family stated for their question in Chapter~\ref{chap:evaluation}.
  \item \emph{Skeleton versus orientations}: a \emph{fixed-graph ablation}.
    The identical per-window graph $M$ feeds allocators that differ only in
    what they read from it: \skelsamp{} symmetrises $M$ before clustering (skeleton
    only), while three orientation-aware allocators (\skelsem{}, \orientsamp{}, \orientsem{};
    Chapter~\ref{chap:sensitivity}) consume the directed matrix itself;
    orientation enters through the allocation step, the ordering step, or
    both.
\end{itemize}

\noindent The research question carries two requirements, each tested in
its own right: the first is the question's opening part, the second a
precondition for its attribution part.
\begin{enumerate}[label=\textbf{Requirement \arabic*.}, itemsep=1pt, topsep=3pt]
  \item \emph{The causal graph must improve on the correlation matrix at
    all}: the question's first part, and the attribution part's
    premise, since without any improvement there is nothing to
    attribute. Tested
    like-for-like: the
    same HRP machinery fed the symmetrised graph (\skelsamp{}) versus the sample
    correlation (\HRP{}), in Section~\ref{sec:r1}.
  \item \emph{An allocator must be able to exploit orientation}. Symmetric
    clustering cannot, so the orientation-aware family of
    Chapter~\ref{chap:sensitivity} is built for the purpose: unit-verified to
    respond to edge direction, while the controls are verified
    orientation-blind under stated invariances (both are transpose-invariant;
    one is invariant to arbitrary per-edge reorientation;
    Chapter~\ref{chap:sensitivity} states which guarantee each carries).
\end{enumerate}

\noindent Granted both, the question's attribution part reduces to two
contrasts with a
shared anchor,
answered in order in Chapter~\ref{chap:evaluation}. The \emph{skeleton
contribution}, $\skelsamp{}-\HRP{}$, measures the value of knowing which
pairs are
linked; the \emph{orientation contribution}, $\orientany{}-\skelsamp{}$,
taken for each of the direction-aware variants in turn, measures the
additional value,
holding the graph fixed, of knowing the direction of each link. Because
both are measured against the same anchor, the two contributions sum
exactly to the total causal-graph gain, so nothing is left unattributed.
Success is therefore any exactly reproducible, rigorously adjudicated answer
for both contributions, including a null: the failure mode the design guards
against is not a small gain but an unattributable one.

A subsidiary question follows once the decomposition is measured: through
what channel does orientation pay, directional content read from the
arrows, or properties of the implied covariance that a graph-free
construction can reproduce? The distinction matters, because only the
first would vindicate a causal reading of the graph.
Chapter~\ref{chap:evaluation} answers it with two covariance controls,
Ledoit--Wolf
shrinkage and a de-factored covariance, and completes the symmetry with a
graph-free skeleton control that asks the same question of the skeleton
channel.

\section{Challenges and contributions}
\label{sec:challenges}

The work is organised as three technical challenges (C1--C3) and the
contributions addressing them (A1--A3), one chapter each. The mapping onto the
research question
is direct: C1 supplies the graph, C2 supplies the allocators that decompose
its value, and C3 supplies the inference that adjudicates the decomposition.
Together they integrate three literatures that have not previously met at
the allocation step (causal structure learning, hierarchical portfolio
construction, and the backtest-rigour programme), and A1--A3 are the
project's objectives: Chapters~\ref{chap:discovery}--\ref{chap:evaluation}
deliver them, and Chapter~\ref{chap:evaluation} assesses them jointly
against the research question.

\begin{enumerate}[label=\textbf{C\arabic*}, itemsep=2pt, topsep=3pt]
  \item \emph{Discovering a credible asset--asset causal graph at every
    rebalance.} The graph must be estimated over a ${\sim}130$-variable joint
    panel of assets and macro drivers, under an economically-motivated
    directional prior (drivers may cause assets; single stocks do not cause
    the 10-year Treasury yield), 215 times per configuration, cheaply enough
    to sweep, and reproducibly: a graph that changes between runs
    invalidates every downstream comparison.
  \item \emph{Feeding a directed graph to an allocator that fundamentally
    consumes symmetric matrices.} Hierarchical clustering requires a
    symmetric distance; naive symmetrisation discards exactly the
    orientations in question. Exploiting them requires allocation machinery
    with no analogue in the hierarchical-allocation literature, and an
    experimental design
    in which orientation is the only treatment variable.
  \item \emph{Attributing small performance differences credibly.} The
    candidate effects are $0.01$--$0.03$ Sharpe on heavy-tailed returns,
    produced by a study that evaluated \rsNtrials\ configurations, so
    the search itself, if unpriced, could manufacture the whole effect.
\end{enumerate}

\begin{enumerate}[label=\textbf{A\arabic*}, itemsep=2pt, topsep=3pt]
  \item \emph{A reproducible rolling pipeline for causal discovery}
    (addresses C1; Chapter~\ref{chap:discovery}): DYNOTEARS with the directional prior
    enforced natively and empirically verified to matter; a content-keyed
    discovery cache that makes the full ablation grid tractable
    (a backtest re-runs in ${\approx}40$ seconds against cached fits
    instead of ${\approx}15$ hours); a data-determinism defect found
    and fixed; and a single extraction chokepoint guaranteeing that every
    allocator receives the identical graph.
  \item \emph{The orientation-aware allocator family and the fixed-graph
    ablation} (addresses C2; Chapter~\ref{chap:sensitivity}), the core
    methodological novelty: a covariance implied by a structural equation
    model (SEM), which propagates shocks through all directed paths of the
    DAG, and
    \emph{causal-ordered bisection}, which replaces the clustering
    dendrogram with the DAG's topological order. Each is deterministic and
    seed-free, and each is evaluated against matched skeleton-only and
    correlation controls on the identical graphs.
  \item \emph{A multiplicity-, distribution- and stability-aware evaluation}
    (addresses C3; Chapter~\ref{chap:evaluation}): stationary-block-bootstrap
    contrasts for every pairwise step of the decomposition; and the
    PSR/DSR/Reality-Check/SPA/MCS battery (the DSR deflating against all
    \rsNtrials\ trials, the family-wise tests over their stated
    per-question families).
\end{enumerate}

\section{Quantified benefits}
\label{sec:benefits}

All benefits below are measured on the protocol of
Section~\ref{sec:question} (net of 5\,bps, monthly rebalances, 2007--2024)
with significance from the Politis--Romano stationary block bootstrap
\cite{politis1994stationary}, and are reported as measured, including
where the answer is small or null. Full detail is in
Chapter~\ref{chap:evaluation}.

\begin{itemize}[itemsep=2pt, topsep=3pt]
  \item \textbf{The decomposition (the answer to the research question).}
    The graph's total gain over the correlation control is stable across
    estimation windows ($+0.018$ to $+0.032$ net Sharpe for \skelsem{}), worth
    $+0.4$ to $+0.7$ percentage points of net compound annual growth rate
    (CAGR) per year, but its composition shifts with the window: the
    skeleton's contribution is hump-shaped and peaks near the conventional
    one-year window, while orientation's is positive at every window and
    largest at the extremes. One pairwise contrast in twelve is
    individually significant, consistent with multiplicity alone, and
    family-wise support is confined to marginal
    rejections under the narrower reported family;
    Chapter~\ref{chap:evaluation} gives the full estimates, the
    anchor-robustness caveats and the multiplicity treatment.
  \item \textbf{What orientation buys is robustness, not level.} The
    orientation-aware allocators are strong at every window (\skelsem{}
    spans $0.395$--$0.415$ and \orientsem{} $0.399$--$0.419$ across the
    four estimation windows, against the skeleton-only allocator's
    $0.372$--$0.403$), and their edge over the skeleton is largest where
    the sample moments feeding the control are weakest (the nine-month and
    two-year windows), consistent with structure substituting for
    sample-moment quality. A direct control identifies the mechanism: the
    SEM covariance's removal of the common market factor, which a
    graph-free de-factored covariance reproduces at every window
    (Chapter~\ref{chap:evaluation}), so the robustness is real but is not
    evidence of directional content. A matching graph-free skeleton
    control shows the converse for the skeleton channel: its one-year
    gain does require the discovered graph. In point estimates the covariance-channel
    edge is invariant to
    transaction costs up to 20\,bps (the ordering route's one-year
    increment erodes across the cost sweep; the sparsity sweep shows no
    threshold sensitivity),
    and at the 18-month window \orientsem{} attains the highest Deflated
    Sharpe ratio of all \rsNtrials\ configurations
    ($\rsDtwosDSRWthree$), with the graph-free de-factored control
    ($\rsDdfDSRWthree$) and \skelsem{} ($\rsDoneDSRWthree$) immediately
    behind it. The effect is, however,
    method-dependent (absent under VARLiNGAM graphs at all four windows)
    and, descriptively, concentrates in calm regimes (the bottom quintile
    of the Cboe Volatility Index, VIX) rather than stress.
  \item \textbf{Novelty relative to the closest published work.} Howard et
    al.\ \cite{howard2025causal} build causal graphs on the S\&P~500 but
    never set weights from them. Setting hierarchical weights
    directly from a causal graph, decomposing skeleton from orientation on
    fixed graphs, and the orientation-aware allocator family itself are, to
    my knowledge, new. Because no published configuration sets hierarchical
    weights from a discovered graph, no external number exists for a direct
    comparison; the external benchmark is therefore the \HRP{} control
    itself, L\'opez de Prado's published textbook construction, evaluated
    on the identical protocol.
\end{itemize}

In summary: \emph{the causal graph's value to hierarchical allocation is
stable across estimation windows, but its source shifts: the skeleton
pays most clearly near the conventional one-year window, while the
orientations pay at every estimation window. Family-wise, support reaches
conventional significance only under the narrower reported family, at the
longest estimation windows, marginally.} The contribution is the allocator family
and the fixed-graph design that made this measurable;
Chapter~\ref{chap:conclusion} restates the full answer with its caveats.

\section{Report structure}
\label{sec:structure}

Chapter~\ref{chap:background} analyses the two works this study builds
on (HRP and the backtest-rigour agenda, and causal graphs in equity
markets) and locates the gap. Chapter~\ref{chap:discovery} (A1) presents the
discovery pipeline that produces the graphs. Chapter~\ref{chap:sensitivity}
(A2) develops the allocator family and the fixed-graph ablation that make the
research question
answerable. Chapter~\ref{chap:evaluation} (A3) evaluates: the two
requirements, the decomposition itself, and how much of it survives rigorous
inference. Chapter~\ref{chap:conclusion} concludes with limitations,
broader impact, and future work. Appendix~\ref{app:tables} collects the full
result tables,
operative hyperparameters and reproducibility notes; the report stands
without it.
